\chapter{Conclusion}
\label{chap:conclusion}

\section{Summary of Findings}
This report proposes and validates a comprehensive Green AI Multi-Agent D3QN model for urban traffic flow. Through strict deterministic testing, we established that a naive pursuit of lower waiting times often triggers network instability. However, integrating chemical tailpipe penalties directly into the Q-learning reward structure coerces the agent into establishing robust, free-flowing corridors that reduce CO\textsubscript{2} and NOx emissions by approximately 25\% while simultaneously curtailing mean delay by 86\%. 

Most critically, the infinitesimal fraction of computing carbon generated confirms that the deployment of sophisticated Deep Learning logic across city infrastructure provides an undeniable net-ecological benefit.

\section{Future Work}
Future research will investigate dynamic tuning mechanisms for the chemical penalty weight ($\alpha_{co2}$) across varying times of day, enabling adaptive transitions between raw throughput maximization and stringent eco-optimization. Furthermore, transferring these decentralized Green AI policies into physical micro-controllers acting on real-world edge hardware constitutes the essential next step toward municipal integration.
